\scanpage{6}
% Two short horizontal strokes at the upper right.
\margin{\mbox{-- --}}

Picard \deleted{\unsure{categories|category}}\inserted{stacks} \unsure{pinned|pinned} by $M.N$, $\sigma$ being the invariant which \inserted{is obtained by\unsure{,}} associating to every section \inserted{$L$} of a Picard stack the symmetry of $\deleted{\unsure{N|L}}L\otimes L$, \unsure{interpreted} as a section of ${}_{2}N$. I know how to prove (unless I am mistaken) that every homomorphism $M\to{}_{2}N$ comes from a suitable Picard stack (pinned by $M,N$) (a priori the obstruction is in $\mathrm{Ext}^{3}(X;M,N)$, but a ``\unsure{universal|universal}'' argument proves that it is zero). This proves that the extension $(*)$ is quite close to being split : \deleted{\unsure{the|a}} every section of the third sheaf, over an arbitrary object of $X$, lifts -- in other words the extension has a ``set-theoretic'' section. Of course, there is in fact better : every section over a $U\in\mathrm{Ob}\,X$ ``comes'' from an element of $\mathscr{H}^{2}(U,{\unsure{\underline{F}|E}}'(M\unsure{.}N\unsure{)|;|,}))$ (hypercohomology $\mathscr{H}^{2}$).

% archive sheet number: 63
